\documentclass[11pt]{article}
\usepackage[margin=1in]{geometry}
\usepackage{amsmath,amssymb}
\usepackage{microtype}

\title{An Analytic Geodesic from a Computational Basis State\\to the Dimension-Four SIC Fiducial}
\author{}
\date{}

\begin{document}
\maketitle

We derive the geodesic from $|0\rangle=(1,0,0,0)^T$ to the analytic
dimension-four SIC fiducial. The target is returned by
\texttt{d4\_sic\_fiducial\_ket()} in \texttt{sics.py}.
We follow the construction in \texttt{utils.py}, implemented by
\texttt{geodesic\_interpolator}.
Throughout, we use the default cyclic Weyl--Heisenberg basis
(\texttt{monomial=False}).

\paragraph{The fiducial.}
Introduce the constants
\begin{equation}
 a=\sqrt{2+\sqrt5},\qquad q=e^{i\pi/4}=\frac{1+i}{\sqrt2},
 \qquad L=\sqrt{2(5+\sqrt5)}.
\end{equation}
The function constructs
\begin{equation}
 |\phi\rangle
 =(H\otimes I_2)\operatorname{diag}(1,\bar q,q,i)
 \frac{(a,1,1,1)^T}{\sqrt{5+\sqrt5}}
 =\frac1L
 \begin{pmatrix}
 a+q\\ \bar q+i\\ a-q\\ \bar q-i
 \end{pmatrix},
 \qquad H=\frac1{\sqrt2}\begin{pmatrix}1&1\\1&-1\end{pmatrix}.
\end{equation}
In particular, its first component is not real. Write
\begin{equation}
 \langle0|\phi\rangle=\frac{a+q}{L}=r e^{i\gamma},
 \quad
 C=|a+q|=\sqrt{a^2+1+\sqrt2\,a},\quad r=\frac CL,
 \quad
 \gamma=\arctan\!\left(\frac1{\sqrt2\,a+1}\right).
\end{equation}

\paragraph{Phase alignment and the orthogonal direction.}
The library first replaces the target by
$|\widetilde\phi\rangle=e^{-i\gamma}|\phi\rangle$, so that
$\langle0|\widetilde\phi\rangle=r>0$. Define
\begin{equation}
 R=\sqrt{L^2-C^2}=\sqrt{7+\sqrt5-\sqrt2\,a},
 \qquad \omega=\arccos r,
 \qquad \sin\omega=\frac RL.
\end{equation}
The normalized component of the aligned target orthogonal to $|0\rangle$ is
\begin{equation}
 |\eta\rangle
 =\frac{|\widetilde\phi\rangle-r|0\rangle}{\sqrt{1-r^2}}
 =\frac{e^{-i\gamma}}R
 \begin{pmatrix}0\\\bar q+i\\a-q\\\bar q-i\end{pmatrix}.
\end{equation}
Thus $\langle0|\eta\rangle=0$ and $\langle\eta|\eta\rangle=1$.

\paragraph{The geodesic used by the library.}
For $0\leq t\leq1$, the normalized horizontal lift of the
Fubini--Study geodesic is
\begin{equation}
 \boxed{
 |g(t)\rangle=\cos(t\omega)|0\rangle+\sin(t\omega)|\eta\rangle
 =\begin{pmatrix}
 \cos(t\omega)\\[2pt]
 \dfrac{e^{-i\gamma}(\bar q+i)}R\sin(t\omega)\\[5pt]
 \dfrac{e^{-i\gamma}(a-q)}R\sin(t\omega)\\[5pt]
 \dfrac{e^{-i\gamma}(\bar q-i)}R\sin(t\omega)
 \end{pmatrix}.}
\end{equation}
This is the analytic expression returned by
\texttt{geodesic\_interpolator} for these endpoints, apart from numerical
roundoff. It satisfies
\begin{equation}
 |g(0)\rangle=|0\rangle,\qquad
 |g(1)\rangle=e^{-i\gamma}|\phi\rangle,\qquad
 \langle g(t)|g(t)\rangle=1.
\end{equation}
The parameter $t\omega$ is linear in $t$, giving constant speed along the
shortest projective geodesic. The curve lies in
$\operatorname{span}\{|0\rangle,|\phi\rangle\}$; its intermediate states
are not required to be SIC fiducials.

\paragraph{Exact ket endpoints.}
To require equality to the supplied ket at $t=1$, including its global
phase, choose the phase-adjusted lift
\begin{equation}
 \boxed{
 |f(t)\rangle=e^{it\gamma}|g(t)\rangle
 =e^{it\gamma}\left[
 \cos(t\omega)|0\rangle+
 \frac{e^{-i\gamma}\sin(t\omega)}R
 \big((\bar q+i)|1\rangle+(a-q)|2\rangle+(\bar q-i)|3\rangle\big)
 \right].}
\end{equation}
It obeys exactly
\begin{equation}
 |f(0)\rangle=|0\rangle,\qquad |f(1)\rangle=|\phi\rangle,
 \qquad \langle f(t)|f(t)\rangle=1.
\end{equation}
Since $|f(t)\rangle\langle f(t)|=|g(t)\rangle\langle g(t)|$,
the two expressions describe the same geodesic of physical states.
The additional phase changes the lift in Hilbert space, not the
Fubini--Study path.

\end{document}
